\section*{Introducción}
El objetivo general de esta tarea es aplicar herramientas de análisis financiero y actuarial para estudiar el riesgo de crédito en bancos, empresas públicas, carteras crediticias y portafolios de bonos. Para ello, se utilizarán modelos como pérdida esperada, capital económico, Z de Altman, modelo de Merton, modelo Logit, CreditMetrics y CVA.

Primero, se analizará la información financiera de distintos bancos mexicanos para observar la evolución de su cartera de crédito, cartera vencida, estimaciones preventivas y capital, con el fin de interpretar su situación financiera y las causas de deterioro o quiebra. Después, se estudiarán diferentes carteras de crédito en dos fechas, calculando medidas de riesgo y comparando los resultados bajo distintos supuestos de correlación.

También se evaluarán cuatro empresas públicas de mercados asiáticos y europeos mediante la Z de Altman y el modelo de Merton, estimando su probabilidad de incumplimiento, pérdida esperada, recuperación y LGD. Posteriormente, estos resultados se usarán para analizar una cartera formada por sus pasivos de largo plazo.

Finalmente, se revisará el modelo Logit de la Circular Única de Bancos, se aplicará CreditMetrics a un portafolio de bonos y se explicarán conceptos relacionados con el riesgo de contraparte, como el CVA y casos reales de pérdidas por mala administración del riesgo crediticio.
